\documentclass[hidelinks]{siamart251216}
\usepackage{amssymb,mathtools}
\usepackage{booktabs,array,graphicx}
\graphicspath{{figures/generated/}}
\usepackage[numbers,sort&compress]{natbib}
\allowdisplaybreaks
\newsiamremark{remark}{Remark}
\headers{Positive-Diagonal Matignon Stability}{M. Lucius}
\title{Exact Positive-Diagonal Matignon Stability\\in Dimension Three}
\author{Maximiliano Lucius}
\begin{document}
\maketitle
\begin{abstract}
We study positive-diagonal multiplicative stability for commensurate fractional-order systems in the Matignon region. For real \(3\times3\) matrices on the full-dimensional strict-\(P(-A)\) stratum, four orbit invariants reduce the complete positive-diagonal orbit to a two-dimensional simplex and fractional D-stability is equivalent to \(\kappa<T_\alpha(\beta)\). The threshold reduces exactly to Cain's classical boundary at integer order; for \(2/3<\alpha<1\), the strict band \(T_1(\beta)<\kappa<T_\alpha(\beta)\) is genuinely fractional and full dimensional. The threshold problem is globally strictly convex in logit coordinates, has a unique smooth optimizer, and admits a sharp classical-order expansion. Dimension three is the smallest real dimension in which the fractional difference from Hurwitz D-stability has nonempty full-dimensional interior. For three-species generalized Lotka--Volterra systems, the invariants become abundance-independent reciprocal-pair and directed-cycle feedback coordinates. High-precision and interval-certified computations independently corroborate the analytic results.
\end{abstract}
\begin{keywords}
D-stability, fractional-order systems, Matignon stability, P-matrices, positive diagonal scaling, generalized Lotka--Volterra systems
\end{keywords}
\begin{MSCcodes}
15A18, 34A08, 93D20
\end{MSCcodes}
\section{Introduction}

Classical D-stability asks whether a real matrix remains stable after multiplication by every positive diagonal matrix. In dimension three the Hurwitz problem has an exact solution due to Cain \citep{Cain1976}, while generalized matrix-stability frameworks extend the spectral target from a half-plane to more general regions \citep{Kushel2019,KushelPavani2022}. For a commensurate Caputo system of order \(0<\alpha<1\), asymptotic stability is instead governed by the Matignon sector
\[
\Sigma_\alpha=\{z\neq0:|\arg z|>\alpha\pi/2\}
\]
\citep{Matignon1996,BrandiburGarrappaKaslik2021}. Fixed-polynomial fractional root criteria and cyclic-network secant conditions are known \citep{CermakNechvatal2017,Bourafa2020,Siami2021}, but they do not explicitly eliminate the universal positive-diagonal multiplier for a generic real \(3\times3\) matrix.

We solve that problem on the full-dimensional strict-\(P(-A)\) stratum. Four positive-left-diagonal orbit invariants,
\[
(\beta_{12},\beta_{13},\beta_{23},\kappa),
\]
reduce the complete positive-diagonal orbit to the open simplex, and the all-diagonal Matignon condition becomes the scalar inequality
\[
\kappa<T_\alpha(\beta).
\]
At \(\alpha=1\), \(T_\alpha\) reduces exactly to Cain's threshold. For \(2/3<\alpha<1\), the strict interval
\[
T_1(\beta)<\kappa<T_\alpha(\beta)
\]
is therefore an exact genuinely fractional D-stable band.

The paper makes four main contributions. First, dimension two is classified exactly and cannot support a full-dimensional genuinely fractional difference class. Second, the real-\(3\times3\) positive-diagonal Matignon problem is reduced to the exact threshold \(T_\alpha\), with a complete characterization of the interior. Third, the threshold objective is globally strictly convex in logit coordinates, giving a unique smooth optimizer, strict monotonicity in \(\alpha\), and a sharp classical-order expansion. These results imply that dimension three is the smallest real dimension with a robust genuinely fractional separation from Hurwitz D-stability. Fourth, for three-species generalized Lotka--Volterra systems the same invariants become abundance-independent reciprocal-pair and directed-cycle feedback coordinates.

The contribution is an explicit low-dimensional solution within the known generalized-D-stability framework, not a new definition of fractional D-stability. Kushel and Pavani \cite{KushelPavani2022} provide an abstract all-diagonal forbidden-boundary framework. Here the diagonal multiplier is eliminated exactly in the generic real three-dimensional Matignon problem. The symmetric slice recovers the single-cycle result of Siami \cite{Siami2021}, rather than replacing it.

Throughout, \emph{genuinely fractional} means stable in the Matignon sense for every positive diagonal multiplier but not Hurwitz D-stable. This orbit-level definition is stronger than exhibiting one matrix that is fractionally stable and integer-order unstable only at \(D=I\).

Section~\ref{sec:framework} gives the orbit reduction, Section~\ref{sec:dimension-two} treats dimension two, Section~\ref{sec:exact-3x3} proves the exact three-dimensional criterion, and Section~\ref{sec:threshold-geometry} develops the convex threshold geometry and minimum-dimension consequence. Section~\ref{sec:ecology} gives the GLV interpretation, Section~\ref{sec:computation} summarizes independent computational certification, and Section~\ref{sec:discussion} positions the result relative to prior work and states its limitations.

\section{Positive-diagonal Matignon stability and orbit reduction}
\label{sec:framework}

For the commensurate Caputo system
\[
{}^C D_t^\alpha x=Ax,
\qquad 0<\alpha<1,
\]
asymptotic stability is equivalent to
\[
\sigma(A)\subset
\Sigma_\alpha
:=
\{z\neq0:|\arg z|>\alpha\pi/2\}
\]
\citep{Matignon1996,BrandiburGarrappaKaslik2021}. Define
\[
\mathcal F_\alpha^{(n)}
=
\{A:\sigma(DA)\subset\Sigma_\alpha
\ \forall D\succ0\text{ diagonal}\},
\]
\[
\mathcal D_H^{(n)}
=
\{A:DA\text{ Hurwitz }\forall D\succ0\},
\qquad
\mathcal P_\alpha^{(n)}
=
\mathcal F_\alpha^{(n)}\setminus\mathcal D_H^{(n)}.
\]
All interiors are taken in the Euclidean topology of \(\mathbb R^{n\times n}\). Because
\[
DA=D(AD)D^{-1},
\]
left and right positive-diagonal multiplication are spectrally equivalent. Our Hurwitz sign convention corresponds to Cain's positive-stability convention after replacing his matrix \(M\) by \(-A\) \citep{Cain1976}.

A real matrix is a P-matrix if every principal minor is positive and a \(P_0\)-matrix if every principal minor is nonnegative. Assume \(-A\) is strict P and, in dimension three, write
\[
p_i=-a_{ii}>0,
\qquad
m_{ij}=\det A[\{i,j\}]>0,
\qquad
q=-\det A>0.
\]
Define the positive-left-diagonal invariants
\[
\beta_{12}=\frac{m_{12}}{p_1p_2},
\quad
\beta_{13}=\frac{m_{13}}{p_1p_3},
\quad
\beta_{23}=\frac{m_{23}}{p_2p_3},
\quad
\kappa=\frac{q}{p_1p_2p_3}.
\]

\begin{lemma}[Positive-diagonal orbit reduction]
\label{lem:simplex-reduction}
Let
\[
D=\operatorname{diag}(d_1,d_2,d_3)\succ0,
\qquad
s=\sum_{i=1}^3p_id_i,
\qquad
x_i=\frac{p_id_i}{s}.
\]
Then \(x\in\Delta_2^\circ=\{x_i>0:\sum_i x_i=1\}\), and
\[
\det(\lambda I-DA)
=
\lambda^3+s\lambda^2+s^2B_\beta(x)\lambda
+s^3\kappa x_1x_2x_3,
\]
where
\[
B_\beta(x)
=
\beta_{12}x_1x_2+
\beta_{13}x_1x_3+
\beta_{23}x_2x_3.
\]
Conversely every \(x\in\Delta_2^\circ\) arises from a positive diagonal \(D\), uniquely up to common positive scale.
\end{lemma}

\begin{proof}
The characteristic coefficients are
\[
s=\sum_i p_id_i,
\qquad
b_D=\sum_{i<j}m_{ij}d_id_j,
\qquad
c_D=qd_1d_2d_3.
\]
Substituting \(d_i=sx_i/p_i\) gives
\[
b_D=s^2B_\beta(x),
\qquad
c_D=s^3\kappa x_1x_2x_3.
\]
Conversely \(d_i=sx_i/p_i\) reconstructs \(D\) for any \(s>0\).
\end{proof}

The radial substitution \(\lambda=sz\) shows that \(s\) changes only eigenvalue magnitudes, not arguments. Thus the full positive-diagonal orbit modulo common scale is exactly the open simplex. The four ratios \((\beta,\kappa)\) are unchanged by positive left-diagonal scaling.

\section{Dimension two: exact obstruction}
\label{sec:dimension-two}

Let
\[
A=\begin{pmatrix}a&b\\c&d\end{pmatrix}.
\]

\begin{theorem}[Exact \(2\times2\) classification]
\label{thm:2x2}
For every \(0<\alpha<1\),
\[
A\in\mathcal F_\alpha^{(2)}
\iff
\det A>0,\qquad a\le0,\qquad d\le0.
\]
\end{theorem}

\begin{proof}
For \(D=\operatorname{diag}(x,y)\succ0\),
\[
\operatorname{tr}(DA)=xa+yd,
\qquad
\det(DA)=xy\det A.
\]
If \(\det A\le0\), either zero is an eigenvalue or \(DA\) has a positive real eigenvalue. If \(a>0\), fixing \(x>0\) and sending \(y\downarrow0\) makes \(DA\) approach a matrix with simple positive eigenvalue \(xa\); continuity gives a positive eigenvalue for small \(y>0\). Hence \(a\le0\), and similarly \(d\le0\).

Conversely, if \(\det A>0\) and \(a,d\le0\), every \(DA\) has positive determinant and nonpositive trace. Its real eigenvalues are therefore negative; a nonreal conjugate pair has nonpositive real part and hence arguments of magnitude at least \(\pi/2>\alpha\pi/2\).
\end{proof}

Classical Hurwitz D-stability imposes the strict trace inequality
\[
xa+yd<0
\qquad
\forall x,y>0,
\]
so under the theorem's conditions it fails exactly when \(a=d=0\). Therefore
\[
\mathcal P_\alpha^{(2)}
=
\left\{
\begin{pmatrix}0&b\\c&0\end{pmatrix}:bc<0
\right\}.
\]

\begin{corollary}[No robust two-dimensional separation]
\label{cor:2x2-empty-interior}
For every \(0<\alpha<1\),
\[
\operatorname{int}_{\mathbb R^4}\mathcal P_\alpha^{(2)}
=
\varnothing.
\]
\end{corollary}

\begin{proof}
The set is contained in the codimension-two subspace \(a=d=0\).
\end{proof}

\section{Exact three-dimensional Matignon threshold}
\label{sec:exact-3x3}

Let \(-A\) be strict P and use the notation of Lemma~\ref{lem:simplex-reduction}. For \(D\succ0\),
\[
\det(\lambda I-DA)
=
\lambda^3+a_D\lambda^2+b_D\lambda+c_D,
\]
with
\[
a_D=s,\qquad
b_D=s^2B_\beta(x),\qquad
c_D=s^3\kappa x_1x_2x_3.
\]

\subsection{Exact cubic boundary}

\begin{lemma}[Exact cubic Matignon boundary]
\label{lem:cubic-boundary}
Let
\[
p(\lambda)=\lambda^3+a\lambda^2+b\lambda+c,
\qquad a,b,c>0,
\]
and \(\pi/3<\theta<\pi/2\). Set
\[
u=\cos\theta,\qquad K=1-4u^2>0,
\]
\[
r_\theta(a,b)
=
\frac{ua+\sqrt{u^2a^2+Kb}}{K},
\]
and
\[
H_\theta(a,b)
=
r_\theta(a,b)^2
\left(a+2u\,r_\theta(a,b)\right).
\]
Then all roots satisfy
\[
|\arg\lambda|>\theta
\]
if and only if
\[
c<H_\theta(a,b).
\]
Moreover,
\[
H_\theta(sa,s^2b)=s^3H_\theta(a,b).
\]
\end{lemma}

\begin{proof}
A positive-coefficient cubic has no nonnegative real root. A root can therefore leave the open angular region only through one of the rays \(\lambda=re^{\pm i\theta}\), \(r>0\). Substituting \(\lambda=re^{i\theta}\) and dividing the imaginary part by \(r\sin\theta\) gives
\[
Kr^2-2uar-b=0,
\]
whose unique positive solution is \(r_\theta(a,b)\). The real part then gives
\[
c=r^2(a+2ur)=H_\theta(a,b).
\]
Thus there is exactly one positive \(c\) at which a ray is touched. For \(c>0\) sufficiently small, one root is a small negative real root and the other two perturb the Hurwitz roots of \(\lambda^2+a\lambda+b\); for \(c\to\infty\), two roots approach the cube-root directions \(\pm\pi/3\), which violate \(\theta>\pi/3\). Hence stability is exactly \(c<H_\theta(a,b)\). Homogeneity follows by \(r_\theta(sa,s^2b)=sr_\theta(a,b)\).
\end{proof}

For \(2/3<\alpha<1\), put
\[
\theta=\frac{\alpha\pi}{2},
\qquad
u=\cos\theta,
\qquad
K=1-4u^2,
\]
and define
\[
r_\alpha(b)
=
\frac{u+\sqrt{u^2+Kb}}{K},
\qquad
h_\alpha(b)
=
r_\alpha(b)^2(1+2u\,r_\alpha(b)).
\]

\subsection{Exact elimination of the diagonal multiplier}

Define
\[
\boxed{
T_\alpha(\beta)
=
\min_{x\in\Delta_2^\circ}
\frac{h_\alpha(B_\beta(x))}
{x_1x_2x_3}.
}
\]
The minimum is attained in the interior because the numerator has a positive continuous extension to the closed simplex while \(x_1x_2x_3\to0\) on its boundary.

\begin{theorem}[Exact real-\(3\times3\) criterion]
\label{thm:C10}
Let \(-A\) be strict P and \(2/3<\alpha<1\). Then
\[
\boxed{
A\in\mathcal F_\alpha^{(3)}
\iff
\kappa<T_\alpha(\beta).
}
\]
\end{theorem}

\begin{proof}
For the simplex coordinate \(x\) associated with \(D\),
\[
\det(\lambda I-DA)
=
\lambda^3+s\lambda^2+s^2B_\beta(x)\lambda
+s^3\kappa x_1x_2x_3.
\]
Lemma~\ref{lem:cubic-boundary} and its homogeneity give
\[
\sigma(DA)\subset\Sigma_\alpha
\iff
\kappa
<
\frac{h_\alpha(B_\beta(x))}
{x_1x_2x_3}.
\]
Lemma~\ref{lem:simplex-reduction} identifies positive diagonal ratios with all \(x\in\Delta_2^\circ\), so the inequality holds for every \(D\succ0\) exactly when \(\kappa<T_\alpha(\beta)\).
\end{proof}

This is an explicit solution of the real three-dimensional Matignon instance of generalized positive-diagonal region stability. The abstract all-\(D\) framework itself is known \citep{Kushel2019,KushelPavani2022}.

\subsection{Interior classification}

\begin{lemma}[\(P_0\) necessity]
\label{lem:P0}
If \(A\in\mathcal F_\alpha^{(3)}\), then
\[
-a_{ii}\ge0,
\qquad
\det A[\{i,j\}]\ge0,
\qquad
-\det A>0.
\]
\end{lemma}

\begin{proof}
If \(a_{ii}>0\), keep the \(i\)-th row scaling fixed and send the other two to zero; the limiting matrix has a positive eigenvalue \(d_i a_{ii}\), which persists for nearby positive scalings. If an order-two principal minor is negative, send the remaining row scaling to zero; the limiting \(2\times2\) principal block has eigenvalues of opposite signs. Both contradict Matignon stability. Finally, at \(D=I\), the real eigenvalue of a real cubic must be negative, while the product of the remaining two eigenvalues is positive; hence \(\det A<0\).
\end{proof}

\begin{corollary}[Strict-P interior]
\label{cor:strictP-interior}
If
\[
A\in\operatorname{int}\mathcal F_\alpha^{(3)},
\]
then \(-A\) is strict P.
\end{corollary}

\begin{proof}
Lemma~\ref{lem:P0} gives nonnegative signed minors. A vanishing order-one minor can be perturbed immediately to the forbidden sign. A singular \(2\times2\) principal block can likewise be perturbed arbitrarily little to negative determinant: use a first-order adjugate perturbation when the block has rank one, and \(\operatorname{diag}(t,-t)\) when it is zero. Thus no such equality can occur at an interior point; the determinant inequality is already strict.
\end{proof}

Kellogg's theorem \citep{Kellogg1972} states that every eigenvalue \(\mu\) of a \(3\times3\) strict P-matrix satisfies
\[
|\arg\mu|<2\pi/3.
\]
Applied to \(-DA\), it yields
\[
|\arg\lambda(DA)|>\pi/3.
\]

\begin{theorem}[Interior of \(\mathcal F_\alpha^{(3)}\)]
\label{thm:interior-F}
For \(0<\alpha\le2/3\),
\[
\operatorname{int}\mathcal F_\alpha^{(3)}
=
\{A:-A\text{ strict P}\}.
\]
For \(2/3<\alpha<1\),
\[
\operatorname{int}\mathcal F_\alpha^{(3)}
=
\{A:-A\text{ strict P},\ \kappa<T_\alpha(\beta)\}.
\]
\end{theorem}

\begin{proof}
For \(\alpha\le2/3\), Kellogg gives sufficiency, Corollary~\ref{cor:strictP-interior} gives necessity, and strict P is open.

For \(\alpha>2/3\), necessity follows from Corollary~\ref{cor:strictP-interior} and Theorem~\ref{thm:C10}. Conversely, near any strict-P matrix the invariants vary continuously and remain in a compact positive box. Since
\[
F_{\alpha,\beta}(x)
\ge
\frac{h_\alpha(0)}{x_1x_2x_3},
\]
all nearby minimizers lie in one compact subset of \(\Delta_2^\circ\); hence \(T_\alpha(\beta)\) is locally continuous. The strict gap \(T_\alpha(\beta)-\kappa>0\) therefore persists under sufficiently small perturbations.
\end{proof}

\subsection{Cain endpoint and the exact fractional band}

As \(\alpha\to1^-\),
\[
h_\alpha(b)\to b,
\]
so
\[
T_\alpha(\beta)\to
T_1(\beta)
=
\min_{x\in\Delta_2^\circ}
\frac{B_\beta(x)}{x_1x_2x_3}.
\]
Since
\[
\frac{B_\beta(x)}{x_1x_2x_3}
=
\frac{\beta_{23}}{x_1}
+\frac{\beta_{13}}{x_2}
+\frac{\beta_{12}}{x_3},
\]
Cauchy--Schwarz gives
\[
\boxed{
T_1(\beta)
=
\left(
\sqrt{\beta_{12}}
+\sqrt{\beta_{13}}
+\sqrt{\beta_{23}}
\right)^2.
}
\]
This is exactly Cain's real-\(3\times3\) D-stability boundary \citep{Cain1976}.

\begin{lemma}[Strict fractional expansion of Cain's boundary]
\label{lem:Talpha-gt-T1}
For \(2/3<\alpha<1\),
\[
T_\alpha(\beta)>T_1(\beta)
\qquad
\forall\beta>0.
\]
\end{lemma}

\begin{proof}
For \(r=r_\alpha(b)\),
\[
b=Kr^2-2ur,
\]
hence
\[
h_\alpha(b)-b
=
4u^2r^2+2ur(r^2+1)>0.
\]
Apply this pointwise to the fractional objective and then minimize.
\end{proof}

\begin{corollary}[Exact genuinely fractional band]
\label{cor:exact-band}
For \(2/3<\alpha<1\),
\[
\boxed{
\operatorname{int}\mathcal P_\alpha^{(3)}
=
\{A:-A\text{ strict P},\
T_1(\beta)<\kappa<T_\alpha(\beta)\}.
}
\]
For \(0<\alpha\le2/3\),
\[
\boxed{
\operatorname{int}\mathcal P_\alpha^{(3)}
=
\{A:-A\text{ strict P},\ \kappa>T_1(\beta)\}.
}
\]
\end{corollary}

\section{Convex geometry and the minimum robust dimension}
\label{sec:threshold-geometry}

Theorem~\ref{thm:C10} reduces the all-diagonal problem to
\[
T_\alpha(\beta)
=
\min_{x\in\Delta_2^\circ}
F_{\alpha,\beta}(x),
\qquad
F_{\alpha,\beta}(x)
=
\frac{h_\alpha(B_\beta(x))}{x_1x_2x_3}.
\]
We now show that this variational problem is globally well behaved.

\subsection{Strict convexity and uniqueness}

Parameterize the simplex by
\[
x_1=\frac{e^{y_1}}{e^{y_1}+e^{y_2}+1},
\qquad
x_2=\frac{e^{y_2}}{e^{y_1}+e^{y_2}+1},
\qquad
x_3=\frac1{e^{y_1}+e^{y_2}+1}.
\]
Set
\[
L(y)=\log(e^{y_1}+e^{y_2}+1),
\]
\[
Q_\beta(y)
=
\log\!\left(
\beta_{12}e^{y_1+y_2}
+\beta_{13}e^{y_1}
+\beta_{23}e^{y_2}
\right),
\]
and
\[
\phi_\alpha(s)=\log h_\alpha(e^s).
\]
Then
\[
G_{\alpha,\beta}(y)
:=
\log F_{\alpha,\beta}(x(y))
=
\phi_\alpha(Q_\beta-2L)+3L-y_1-y_2.
\]

\begin{lemma}[Elasticity bounds]
\label{lem:elasticity-sub}
For \(2/3<\alpha<1\) and \(b>0\),
\[
0<
E_\alpha(b)
:=
\frac{b h_\alpha'(b)}{h_\alpha(b)}
<
\frac32,
\qquad
E_\alpha'(b)>0.
\]
Equivalently,
\[
0<\phi_\alpha'(s)<3/2,
\qquad
\phi_\alpha''(s)>0.
\]
\end{lemma}

\begin{proof}
With
\[
u=\cos(\alpha\pi/2),
\qquad
K=1-4u^2,
\]
parameterize the cubic boundary by
\[
b=r(Kr-2u),
\qquad
h=r^2(1+2ur),
\qquad
r>2u/K.
\]
Then
\[
E
=
\frac{b}{h}\frac{dh/dr}{db/dr}
=
\frac{(Kr-2u)(1+3ur)}
{(Kr-u)(1+2ur)}.
\]
The two factors
\[
\frac{Kr-2u}{Kr-u},
\qquad
\frac{1+3ur}{1+2ur}
\]
increase strictly with \(r\), while \(E\to0\) as \(r\downarrow2u/K\) and \(E\to3/2\) as \(r\to\infty\). Since \(b(r)\) is strictly increasing, the result follows.
\end{proof}

\begin{theorem}[Global strict logit convexity]
\label{thm:strict-convexity}
For every \(\beta>0\) and \(2/3<\alpha<1\), \(G_{\alpha,\beta}\) is globally strictly convex and coercive. Hence \(F_{\alpha,\beta}\) has a unique nondegenerate minimizer
\[
x^*(\alpha,\beta)\in\Delta_2^\circ,
\]
and \(T_\alpha(\beta)\) and \(x^*(\alpha,\beta)\) depend smoothly on \((\alpha,\beta)\).
\end{theorem}

\begin{proof}
Writing \(s=Q_\beta-2L\),
\[
\nabla^2G
=
\phi_\alpha''(s)\nabla s\nabla s^\top
+
\phi_\alpha'(s)\nabla^2Q_\beta
+
\bigl(3-2\phi_\alpha'(s)\bigr)\nabla^2L.
\]
Here \(\nabla^2Q_\beta\succeq0\), while \(\nabla^2L\succ0\) because \(L\) is the log-sum-exp of the three affinely independent vectors \((1,0),(0,1),(0,0)\) with positive softmax weights. Lemma~\ref{lem:elasticity-sub} makes every scalar coefficient nonnegative and the coefficient of \(\nabla^2L\) strictly positive. Thus
\[
\nabla^2G(y)\succ0
\qquad
\forall y\in\mathbb R^2.
\]
As \(\|y\|\to\infty\), \(x(y)\) approaches the simplex boundary, \(x_1x_2x_3\to0\), and the numerator remains strictly positive; hence \(G\to\infty\). Coercivity gives a unique global minimizer, positive definiteness gives nondegeneracy, and the implicit-function theorem yields smooth dependence.
\end{proof}

\subsection{Symmetric slice and canonical family}

\begin{corollary}[Fully symmetric slice]
\label{cor:symmetric-slice}
If
\[
\beta_{12}=\beta_{13}=\beta_{23}=b>0,
\]
then
\[
x^*=(1/3,1/3,1/3),
\qquad
T_\alpha(b,b,b)=27h_\alpha(b/3).
\]
For \(b=1\),
\[
T_\alpha(1,1,1)=1+R_3(\alpha)^3,
\qquad
R_3(\alpha)
=
\frac{\sin(\alpha\pi/2)}
{\sin(\alpha\pi/2-\pi/3)}.
\]
\end{corollary}

\begin{proof}
Permutation symmetry and uniqueness force \(x^*\) to be the simplex center. Substitution gives the first identity; the second simplifies to the \(n=3\) single-cycle threshold of Siami \cite{Siami2021}.
\end{proof}

Consider
\[
A_\gamma
=
\begin{pmatrix}
-1&0&-\gamma\\
\gamma&-1&0\\
0&\gamma&-1
\end{pmatrix},
\qquad
\gamma>0.
\]
It has
\[
\beta=(1,1,1),
\qquad
\kappa=1+\gamma^3.
\]

\begin{corollary}[Exact cyclic threshold]
\label{cor:cyclic-family}
For \(2/3<\alpha<1\),
\[
A_\gamma\in\mathcal F_\alpha^{(3)}
\iff
\gamma<R_3(\alpha),
\]
whereas
\[
A_\gamma\in\mathcal D_H^{(3)}
\iff
\gamma<2.
\]
Thus
\[
2<\gamma<R_3(\alpha)
\]
is an exact one-parameter genuinely fractional D-stable band.
\end{corollary}

\begin{proof}
The fractional statement follows from Theorem~\ref{thm:C10} and Corollary~\ref{cor:symmetric-slice}; the classical endpoint is
\[
1+\gamma^3<T_1(1,1,1)=9.
\]
\end{proof}

\subsection{Order dependence and the classical limit}

\begin{theorem}[Strict monotonicity]
\label{thm:alpha-monotonicity}
For fixed \(\beta>0\),
\[
\frac23<\alpha_1<\alpha_2\le1
\quad\Longrightarrow\quad
T_{\alpha_1}(\beta)>T_{\alpha_2}(\beta).
\]
\end{theorem}

\begin{proof}
For fixed \(a,b>0\), the cubic boundary value \(H_\theta(a,b)\) strictly decreases as the forbidden angle \(\theta\) increases: a cubic touching the larger-angle ray lies strictly inside the wider Matignon region associated with the smaller angle. Hence
\[
h_{\alpha_1}(b)>h_{\alpha_2}(b)
\qquad
\forall b>0.
\]
Evaluating the two objectives at a minimizer of the \(\alpha_1\) problem gives the stated strict inequality.
\end{proof}

\begin{figure}[t]
\centering
\includegraphics[width=.68\linewidth]{fig_threshold_ratio.pdf}
\caption{Fractional enlargement of Cain's threshold for representative invariant triples. Theorem~\ref{thm:alpha-monotonicity} gives strict decrease in \(\alpha\), and Theorem~\ref{thm:C14} gives convergence to \(T_1\) at integer order. The ordinate is logarithmic.}
\label{fig:threshold-ratio}
\end{figure}

Set
\[
S=\sqrt{\beta_{12}}+\sqrt{\beta_{13}}+\sqrt{\beta_{23}},
\qquad
\Gamma=\sqrt{\beta_{12}\beta_{13}\beta_{23}}.
\]

\begin{theorem}[Sharp classical-limit rate]
\label{thm:C14}
As \(\alpha\to1^-\),
\[
T_\alpha(\beta)
=
T_1(\beta)
+
C(\beta)(1-\alpha)
+
O((1-\alpha)^2),
\]
where
\[
T_1(\beta)=S^2,
\qquad
C(\beta)
=
\pi
\left(
\frac{S^{5/2}}{\sqrt\Gamma}
+
S^{3/2}\sqrt\Gamma
\right).
\]
\end{theorem}

\begin{proof}
Put \(\varepsilon=1-\alpha\). Since
\[
u=\cos(\alpha\pi/2)
=
\frac{\pi}{2}\varepsilon+O(\varepsilon^3),
\]
implicit differentiation of the normalized boundary equation at \(u=0\) gives
\[
h_\alpha(b)
=
b+\pi\varepsilon\sqrt b(1+b)+O(\varepsilon^2).
\]
At \(\alpha=1\) the unique simplex minimizer is
\[
x_1^*=\frac{\sqrt{\beta_{23}}}{S},
\quad
x_2^*=\frac{\sqrt{\beta_{13}}}{S},
\quad
x_3^*=\frac{\sqrt{\beta_{12}}}{S},
\]
with
\[
x_1^*x_2^*x_3^*=\frac{\Gamma}{S^3},
\qquad
B_\beta(x^*)=\frac{\Gamma}{S}.
\]
The envelope theorem, justified by Theorem~\ref{thm:strict-convexity}, therefore gives
\[
C(\beta)
=
\pi
\frac{\sqrt{\Gamma/S}\,(1+\Gamma/S)}
{\Gamma/S^3},
\]
which simplifies to the stated formula.
\end{proof}

\subsection{Minimum dimension for robust genuinely fractional separation}

\begin{theorem}[Minimum robust dimension]
\label{thm:C09}
For every \(0<\alpha<1\),
\[
\min\left\{
n:
\operatorname{int}
\left(
\mathcal F_\alpha^{(n)}
\setminus
\mathcal D_H^{(n)}
\right)
\neq\varnothing
\right\}
=
3.
\]
\end{theorem}

\begin{proof}
For \(n=1\), the two stability notions coincide. For \(n=2\), Corollary~\ref{cor:2x2-empty-interior} gives empty interior.

For \(n=3\), use \(A_\gamma\). It satisfies \(-A_\gamma\) strict P for every \(\gamma>0\), and it is Hurwitz D-stable exactly for \(\gamma<2\). If \(0<\alpha\le2/3\), Kellogg's theorem gives \(A_\gamma\in\mathcal F_\alpha^{(3)}\) for every \(\gamma\), so any \(\gamma>2\) lies in the genuinely fractional class. If \(2/3<\alpha<1\), Corollary~\ref{cor:cyclic-family} gives fractional D-stability for \(\gamma<R_3(\alpha)\), and \(R_3(\alpha)>2\) because
\[
T_\alpha(1,1,1)>T_1(1,1,1)=9.
\]
Choose \(2<\gamma<R_3(\alpha)\). In both regimes the defining inequalities are strict, so the difference class contains an open neighborhood of the chosen matrix.
\end{proof}

The proof-complete technical version of the manuscript also develops a global parametrization of the threshold surface, exact invariant-space sensitivities, the \(\alpha\downarrow2/3\) asymptotic, a closed-form sufficient certificate, and a general realizability theorem. Those secondary results are retained in the accompanying technical manuscript and reproducibility archive but are not needed for the central elimination theorem or the minimum-dimension conclusion.

\section{Three-species ecological interpretation}
\label{sec:ecology}

Consider the commensurate fractional GLV system
\[
{}^C D_t^\alpha x_i
=
x_i\left(r_i+\sum_{j=1}^3a_{ij}x_j\right),
\]
a positive equilibrium \(x^*\gg0\) has Jacobian
\[
J=\operatorname{diag}(x^*)A.
\]
The results below concern local Jacobian stability. Principal-minor and cycle decompositions are classical in ecological stability analysis \cite{ClarkHallam1982,Kinkhabwala2015}. Here the exact fractional all-diagonal boundary is written directly in those coordinates.

\begin{proposition}[Positive-equilibrium orbit equivalence]
\label{prop:GLV-invariance}
For every \(x^*\gg0\),
\[
J\in\mathcal F_\alpha^{(3)}
\iff
A\in\mathcal F_\alpha^{(3)},
\qquad
J\in\mathcal D_H^{(3)}
\iff
A\in\mathcal D_H^{(3)}.
\]
Hence \(J\in\mathcal P_\alpha^{(3)}\iff A\in\mathcal P_\alpha^{(3)}\), and the four invariants \((\beta,\kappa)\) are unchanged.
\end{proposition}

\begin{proof}
If \(E=\operatorname{diag}(x^*)\), then
\[
\{DJ:D\succ0\}
=
\{DEA:D\succ0\}
=
\{\widetilde D A:\widetilde D\succ0\},
\]
because \(D\mapsto DE\) is a bijection of the positive diagonal cone. The invariant ratios are unchanged by the same row scaling.
\end{proof}

Write
\[
p_i=-a_{ii}>0,
\qquad
g_{ij}=\frac{a_{ij}a_{ji}}{p_ip_j}.
\]
Then
\[
\boxed{\beta_{ij}=1-g_{ij}.}
\]
Thus opposite-sign reciprocal interactions have \(g_{ij}<0\) and increase \(\beta_{ij}\), whereas same-sign reciprocal interactions have \(g_{ij}>0\) and decrease it, subject to strict P.

For the two oriented three-cycle products,
\[
\ell_{123}
=
\frac{a_{12}a_{23}a_{31}}{p_1p_2p_3},
\qquad
\ell_{132}
=
\frac{a_{13}a_{32}a_{21}}{p_1p_2p_3},
\]
put
\[
L_3=\ell_{123}+\ell_{132}.
\]

\begin{proposition}[Determinant-loop identity]
\label{prop:loop-identity}
\[
\boxed{
\kappa
=
\beta_{12}+\beta_{13}+\beta_{23}-2-L_3.
}
\]
\end{proposition}

\begin{proof}
Expanding \(\det A\) and dividing by \(p_1p_2p_3\) gives
\[
\frac{\det A}{p_1p_2p_3}
=
-1+L_3+g_{12}+g_{13}+g_{23}.
\]
Use \(\kappa=-\det A/(p_1p_2p_3)\) and \(g_{ij}=1-\beta_{ij}\).
\end{proof}

\begin{corollary}[Exact loop-space fractional band]
\label{cor:loop-band}
For \(2/3<\alpha<1\), set
\[
C_\beta=\beta_{12}+\beta_{13}+\beta_{23}-2.
\]
On the strict-P stratum,
\[
A\in\operatorname{int}\mathcal P_\alpha^{(3)}
\]
if and only if
\[
\boxed{
C_\beta-T_\alpha(\beta)
<
L_3
<
C_\beta-T_1(\beta).
}
\]
\end{corollary}

\begin{proof}
Substitute \(\kappa=C_\beta-L_3\) into Corollary~\ref{cor:exact-band}.
\end{proof}

\begin{figure}[t]
\centering
\includegraphics[width=.68\linewidth]{fig_ecological_band.pdf}
\caption{Cain and fractional loop-space boundaries for \(\beta=(3,0.5,1)\). The genuinely fractional region lies between the two curves.}
\label{fig:ecological-band}
\end{figure}

The ecological interpretation is therefore invariant under positive equilibrium abundance scaling. It separates normalized reciprocal-pair feedback \(\beta_{ij}=1-g_{ij}\) from total directed three-cycle feedback \(L_3\). This is a structural local-stability statement; it does not assert equilibrium feasibility, global persistence, or a causal effect of changing one biological interaction coefficient in isolation.

\section{Computational certification and reproducibility}
\label{sec:computation}

All theorem claims above are analytic. Computation was used only as an independent falsification and certification layer. Two routes were kept separate: an invariant-space implementation of
\[
T_\alpha(\beta)
=
\min_{x\in\Delta_2^\circ}
\frac{h_\alpha(B_\beta(x))}{x_1x_2x_3},
\]
and direct optimization of the spectral margin
\[
\min_{D\succ0}
\left[
\min_{\lambda\in\sigma(DA)}
|\arg\lambda|-\frac{\alpha\pi}{2}
\right].
\]
The latter uses only eigenvalues of the full matrix \(DA\), so agreement tests the elimination theorem rather than a second implementation of the same formula.

Across two campaigns, \(2{,}180{,}000\) adversarial matrix/order cases produced no genuine discrepancy with Theorem~\ref{thm:C10}; all \(56{,}234\) numerically delicate cases were independently recomputed in multiprecision. A second atlas enclosed \(938\) threshold values by interval arithmetic. Of these, \(903\) had a global branch-and-bound certificate independent of Theorem~\ref{thm:strict-convexity}; \(35\) extreme-anisotropy cases used the proved convexity/uniqueness result to complete localization. The worst relative interval width was \(1.1\times10^{-35}\), and every independent 50-digit optimizer lay inside its reported enclosure.

The \(35\)-matrix canonical library includes symmetric and nonsymmetric examples, matrices near the Cain and fractional boundaries, mixed-sign ecological examples, and matrices that are Hurwitz at \(D=I\) but not Hurwitz D-stable. One representative is
\[
A=
\begin{pmatrix}
-1&-3/2&0\\
-1/2&-1&-3\\
-3&0&-2
\end{pmatrix},
\qquad
\alpha=0.7.
\]
Here
\[
\beta=(1/4,1,1),
\qquad
\kappa=7,
\qquad
T_1=6.25,
\qquad
T_\alpha\approx4726.94.
\]
Although
\[
\max\Re\sigma(A)\approx-0.0789,
\]
the Cain-minimizing positive diagonal gives
\[
\max\Re\sigma(D_HA)\approx0.0384>0.
\]
Thus \(A\) itself is Hurwitz, its classical positive-diagonal orbit is not, but its full orbit is Matignon-stable at \(\alpha=0.7\).

The C-15 differential structure was also checked at \(8{,}000\) parameter points. The \(500\) worst floating-point cases were recomputed at 40 digits; the identity
\[
\frac{\partial T_\alpha}{\partial\beta_{ij}}
=
\frac{h_\alpha'(B^*)}{x_k^*}
\]
then agreed to at worst \(2.4\times10^{-33}\), and the logit Hessian remained positive definite in every tested case. The final repository test suite has \(135\) passing tests and no failures.

All scripts, fixed seeds, software metadata, interval certificates, canonical matrices, and figure-source data are maintained in the public repository
\[
\text{\url{https://github.com/maximilianolucius/fractional-d-stability-networks}}.
\]
The numerical evidence is reproducible but logically secondary: no theorem is inferred from finite sampling.

\section{Discussion}
\label{sec:discussion}

The main result is an explicit solution of a known generalized-D-stability problem in the first dimension where a robust genuinely fractional separation is possible. Generalized multiplicative region stability is not new \citep{Kushel2019}; \cite{KushelPavani2022} give an abstract all-positive-diagonal forbidden-boundary framework. Our contribution is the exact elimination of that diagonal quantifier for the generic real \(3\times3\) Matignon problem on the full-dimensional strict-\(P(-A)\) stratum.

The closest classical predecessor is Cain's complete real-\(3\times3\) Hurwitz D-stability criterion \citep{Cain1976}. Here it appears exactly as the endpoint
\[
T_1(\beta)
=
\left(
\sqrt{\beta_{12}}+
\sqrt{\beta_{13}}+
\sqrt{\beta_{23}}
\right)^2,
\]
while Theorem~\ref{thm:alpha-monotonicity} shows that \(T_\alpha>T_1\) for every fractional order in the high-order regime. Bahl--Cain fixed-inertia theory \citep{BahlCain1977} is complementary but cannot determine Matignon stability because the fractional stable region includes a right-half-plane sliver and therefore requires angular information finer than inertia. The cyclic secant theorem of \cite{Siami2021} is recovered exactly as the symmetric invariant slice.

Classical sufficient computational tests, including LMI approaches \citep{OliveiraPeres2004}, and scalable determinant and principal-minor conditions \citep{Kushel2023,Kushel2026} address a different tradeoff: tractable sufficient testing in general dimension rather than an exact angular elimination in dimension three.

Theorem~\ref{thm:C09} explains why \(n=3\) is not merely a convenient small example. In dimension two the fractional-only class is lower-dimensional; in dimension three it has nonempty interior for every \(0<\alpha<1\). The positive-diagonal orbit reduction extends to higher dimension, but from \(n=4\) onward the normalized polynomial contains several independent angular coefficients and the exact scalar threshold structure is lost. Exploratory \(4\times4\) computations are therefore reserved for future work.

The strict-P hypothesis should also be interpreted correctly. Corollary~\ref{cor:strictP-interior} shows that every full-dimensional interior point lies on that stratum, so the theorem completely describes the robust class. It does not classify every lower-dimensional boundary component with vanishing signed principal minors. This distinction parallels classical work on the interior and strong D-stability \citep{Hartfiel1980,Abed1986}.

Finally, the GLV interpretation is local. Positive equilibrium abundances reparameterize the same positive-diagonal orbit, and the loop variables give invariant coordinates for that orbit. They do not by themselves imply feasibility, global persistence, nonlinear basin stability, or a one-parameter causal intervention rule.

The principal limitations are therefore explicit: common fractional order, real matrices, full-dimensional classification, dimension three, and local equilibrium stability. These boundaries define the natural continuation of the theory rather than qualifications of the exact three-dimensional result.

\section{Conclusion}

For real \(3\times3\) matrices on the full-dimensional strict-\(P(-A)\) stratum, positive-diagonal Matignon stability is exactly equivalent to
\[
\kappa<T_\alpha(\beta),
\]
where four invariant principal-minor ratios reduce the entire positive-diagonal orbit to a two-dimensional simplex. The threshold recovers Cain's classical D-stability surface at \(\alpha=1\), is globally strictly convex after a logit transformation, and has a unique smooth optimizer.

The strict band
\[
T_1(\beta)<\kappa<T_\alpha(\beta)
\]
therefore gives an exact orbit-level separation between fractional and Hurwitz D-stability. Together with the complete \(2\times2\) classification and the explicit cyclic family, this shows that dimension three is the smallest real dimension in which the difference has nonempty full-dimensional interior for every \(0<\alpha<1\).

For three-species generalized Lotka--Volterra systems, positive equilibrium abundance scaling does not change the positive-diagonal orbit. The exact threshold can therefore be expressed in abundance-independent reciprocal-pair and directed-cycle feedback coordinates. Independent high-precision and interval computations found no discrepancy with the analytic theory.

The general simplex reduction persists beyond dimension three, but the higher-order angular geometry does not collapse to the same scalar cubic boundary. Extending exact positive-diagonal Matignon stability to structured \(4\times4\) families is the next natural problem.

\section*{Declaration on generative AI use}

Generative AI tools, including OpenAI ChatGPT and Anthropic Claude, were used to assist with drafting and editing, organization of literature searches, mathematical proof exploration and consistency checking, development and review of symbolic and numerical code, and preparation of reproducibility materials. These tools are not authors. The author assumes responsibility for all content.

\bibliographystyle{siamplain}
\bibliography{references}
\end{document}
